% =====================================================================================
% theory_section.tex -- main-text theory section, for insertion after Background (Section 2).
% Written 2026-09-25, after the results in Section 4 were seen; revised the same day after an adversarial review
% (see THEORY_NOTES.md section 10). New references: paper/refs_theory_pending.bib (Crossref-verified); all other
% keys are in refs_merged.bib. Uses \ci (defined in main.tex). Cross-reference used: sec:results (main.tex).
% \SItheory names the SI theory appendix (paper/theory_si.tex, label si:theory); on insertion, redefine it as
% \SI{<n>} with the final SI section number.
\providecommand{\SItheory}{the SI theory appendix}
% =====================================================================================
\section{Two audiences, one name: a model of coupling}\label{sec:theory}

Why should faculty hiring measure academic standing, and program pay the labour market's regard for
graduates? This section derives both readings from one model, states the assumptions under which each holds and
what breaks them, and says what $\rho_f$ then measures. We wrote the model after the results of
Section~\ref{sec:results} were known. It organises them after the fact, and none of them is evidence for it. Its
out-of-sample content is a set of registered tests that have not been run, listed at the end of this section.
The proofs, and a sorting of every result by how the model accommodates it, are in \SItheory.

\paragraph{Two audiences, two products, one name.} A university $i$ sells two products of its department in
field $f$ under one name: doctoral graduates, bought by the field's hiring committees, and bachelor's graduates,
bought by employers. Following \citet{macleod2015reputation} and \citet{tadelis1999name}, a unit's
\emph{reputation} with an audience is that audience's expectation of the quality of what it buys from the unit,
given what it observes. It is a belief held by a named population about a named product. Beliefs are not
observed, but repeated costly choices reveal the orderings they imply, much as matriculation decisions reveal
applicants' ranking of colleges \citep{avery2013revealed}. We call such an ordering the audience's
\emph{valuation}. Write the academy's valuation of department $(i,f)$ as $s_{if}=\alpha_fa_i+r_{if}$ and the mean
productivity of its bachelor's graduates, selection included, as $\theta_{if}=\mu_f+\beta_fx_i+y_{if}$. Here
$a_i$ and $x_i$ are institution components, $r_{if}$ and $y_{if}$ are departmental deviations, and $\beta_f$ is
how strongly field $f$'s labour market rewards the institution component. The two sides are correlated within a
level ($\mathrm{Corr}(a,x)=\chi$, $\mathrm{Corr}(r,y)=\kappa_f$). We assume that they are uncorrelated across
levels, for example $\mathrm{Cov}(r,x)=0$; departments that are relatively strong at highly selective
institutions would break this. Employers see the name and the major on a CV and pay their posterior expectation
of a graduate's productivity, given what they have seen of the unit's past graduates. Attribution is costly and
labour markets are segmented, so what counts is how many graduates of a unit a given employer has observed and
attributed to it. With $n_u$ such graduates, normal learning puts weight
$\omega_u=n_u\tau_u^2/(n_u\tau_u^2+\sigma^2)$ on the unit's own record ($\tau_u^2$ is the prior variance of unit
means, $\sigma^2$ that of individual performance). A name pools all of an institution's graduates, so
$\omega_{\text{name}}\approx1$. A department's record counts only where employers attribute performance to
departments. The pooling is analogous to umbrella branding \citep{wernerfelt1988umbrella}, read as buyers'
inference across the products of one brand \citep{choi1998brand,cabral2000stretching} rather than as a firm's
signalling choice.

\paragraph{Proposition 1 (what the indicators reveal).} \emph{(a) Hiring.} Let $A_{uv}$ count field-$f$ faculty
at $v$ with a PhD from $u$. Assume vertical sorting net of volume: in every pair of departments that exchanges
faculty, more hires are expected to flow from the higher-valued to the lower-valued department, and
production-to-hiring ratios do not set that direction. Then the ordering that minimises upward hires agrees with
the academy's valuation on every exchanging pair, and on every pair linked by a chain of exchanging pairs that is
monotone in the valuation; pairs linked by no such chain are not ordered by the data. The sample version
converges as hires accumulate. This is a consistency property. The substantive assumption is vertical sorting
net of volume, which is the construct's definition and can fail. Two gaps separate it from the published $F$.
First, $F$ is SpringRank, not the minimum-violation ordering. Under SpringRank's generative model,
$\mathbb E A_{uv}=c\exp\{-\tfrac{\beta}{2}(s_u-s_v-1)^2\}$ with a single constant $c$, the direction of exchange
given the pair total is logistic in $2\beta(s_u-s_v)$, a Bradley--Terry comparison \citep{bradley1952rank}, and
SpringRank approaches the maximum-likelihood ranks when hierarchy is strong \citep{debacco2018physical}. Under
vertical sorting alone the two orderings can disagree (\SItheory). Second, when expected flows scale with
production and hiring, $\mathbb E A_{uv}=p_uh_v\varphi(s_u-s_v)$, the direction of exchange also depends on
$(p_u/h_u)/(p_v/h_v)$. SpringRank has no degree correction, so at equal valuation it ranks net exporters higher,
and faculty production is highly unequal \citep{wapman2022quantifying}. $F$ is therefore a department's
net-placement position. It is the academy's valuation only if production-to-hiring ratios do not set the
direction of flows (test T0 below).

\emph{(b) Pay.} Suppose employers compete for graduates and pay the posterior expectation of productivity, see
institution and major at hiring, and set pay freely rather than by a schedule, and that program medians cover the
graduates who work for them. Then, with the name's mean known to employers ($\omega_{\text{name}}\approx1$; the
general form is in \SItheory), the program's log median $t$ years after graduation is
\begin{equation}\label{eq:theory-wage}
m_{if}(t)=\mu_f(t)+\beta_fx_i+\Lambda_{ft}\,y_{if}+g_{i}(t),\qquad \Lambda_{ft}=\omega_f+(1-\omega_f)\lambda_t .
\end{equation}
Here $\lambda_t=b+(1-b)\,t/(t+k)$ is the weight employers put on a graduate's own revealed performance ($b$ for
what they see at hiring, $k$ for the speed of learning), and $g_i(t)$ is growth that differs across
institutions, for example through access to firms with internal promotion ladders.

\paragraph{What the two readings license.} Part (a) warrants reading $F$ as \emph{academic prestige}: the
ordering of departments revealed by where their PhDs are placed, which sociologists call status in a deference
network \citep{burris2004academic,podolny1993status,sauder2012status}. PhD origin shapes placement beyond
productivity \citep{long1979entrance}, so this status may partly reinforce itself. The analogy with revealed
preference is loose. No one chooses between $u$ and $v$; faculty hiring is two-sided matching, shaped by
candidates' location and dual-career constraints; and PhDs who are not placed are not observed. The reading fails
where subfield, geography or dual careers steer hiring in ways correlated with standing, and where volume sets
the direction of flows. Reading $F$ further as a belief about the quality of a department's PhDs requires one
index to govern both where its PhDs are hired and whom it hires (test T1). Part (b) warrants reading program pay
as the \emph{labour-market valuation} of a program's graduates. At the name, $\omega\approx1$, so employers'
belief and the realised mean productivity of the institution's graduates coincide, selection included. Much of
the premium of selective colleges reflects whom they admit \citep{dale2002estimating,dale2014estimating}, and
their causal effects concentrate in top-tail outcomes that medians do not register
\citep{zimmerman2019elite,chetty2023diversifying}. We reserve \emph{employer reputation} for what the design of
\citet{macleod2017big} identifies: a school-mean premium with graduates' own attainment held fixed. The UK data
identify it coarsely, since attainment bands are observed there: 66--90\% of coupling survives the attainment
control, and within a band selectivity keeps $+0.252$ beyond $G$ (test T4 registers how it changes with
experience). Three further limits apply.
\begin{itemize}
\item Raw pay contains local price levels, which lie outside the construct, and access to high-wage labour
markets. No result here nets out where graduates work (test T7), so every construct claim about pay below is
``not netted of destination location''.
\item Placement into better-paid sectors counts as the market's valuation only if employers ration access to
those sectors. Graduates' preferences, family background and peers can produce the same placement (test T8).
\item Where pay follows a schedule, the wage stops revealing valuation (Proposition~5).
\end{itemize}

\paragraph{Proposition 2 (what coupling measures).} Under Proposition~1, with $F=s+e_F$ and zero covariances
across levels,
\begin{equation}\label{eq:theory-decomp}
\mathrm{Cov}\{F_{if},m_{if}(t)\}=\underbrace{\alpha_f\beta_f\,\chi\,\sigma_a\sigma_x}_{\text{name}}
+\underbrace{\Lambda_{ft}\,\kappa_f\,\sigma_r\sigma_y}_{\text{department}}
+\underbrace{\alpha_f\,\mathrm{Cov}\{a_i,g_{i}(t)\}}_{\text{growth}} .
\end{equation}
$G$ pools the hires of all fields, so it carries only the name and growth terms. The partial covariance of $F$
with pay given $G$ isolates the department term only if the cross-level covariances are zero. For Spearman
correlations computed on the same institutions, the identity $c_F=r_{FG}\,c_G+\sqrt{1-r_{FG}^2}\;\mathrm{sr}_F$
is exact. Here $c_X$ is the coupling of $X$ with pay, $r_{FG}$ the correlation of $F$ with $G$, and
$\mathrm{sr}_F$ the semi-partial correlation of pay with the part of $F$ orthogonal to $G$. $F$ out-predicts
$G$ only if $\mathrm{sr}_F>c_G\tan(\vartheta/2)$, where $\cos\vartheta=r_{FG}$. Three readings follow.
\begin{itemize}
\item Coupling measures agreement between two audiences' orderings of the same institutions, not agreement on
what they value. An academy that values research and a market that values graduates' ability agree whenever the
two covary across institutions ($\chi>0$). That they do is a primitive of the model, not a result.
\item If one institution factor dominates, each field's coupling is proportional to its $\beta_f$, and the field
map is a map of how each field's labour market prices institution-wide status. It would look the same with any
status index, so its interpretation does not depend on $F$ being academic prestige. The construct reading of $F$
matters for the claims below the name. Test T5 checks the invariance.
\item Average SAT and the admission rate are noisy measures of the selection part of $x_i$, the ability of an
institution's intake \citep{macleod2015reputation,macleod2017big}. Partialling them out removes the part of the
market's valuation that is measured intake ability. What remains bundles value added, peers, networks, growth
and unmeasured ability. For the question ``does the academy's standing carry information that the market prices
beyond intake?'', selectivity is a confound. SAT and admission rate alone take mean coupling from $+0.43$ to
$+0.14$ (47 fields). Pell share, control and the state earnings level alone take it to $+0.31$, so raw $\rho_f$
also contains location and composition, which lie outside both constructs.
\end{itemize}

\paragraph{Proposition 3 (below the name).} The department term of \eqref{eq:theory-decomp} depends on when
pay is observed. At entry, if employers see nothing of a graduate beyond the credential ($b=0$), it equals
$\omega_f\kappa_f\sigma_r\sigma_y$. It is then non-zero only if department standing carries information about
graduates' productivity ($\kappa_f\neq0$) and employers attribute performance to departments ($\omega_f>0$).
After entry it needs only $\kappa_f\neq0$, because $\Lambda_{ft}\ge\lambda_t$: each graduate's own record prices
the department deviation whether or not anyone tracks departments. Employer learning is fast.
\citet{lange2007speed} put $k$ near 3, which gives $\lambda_4\approx0.57$, and for college graduates ability is
largely revealed at hiring \citep{arcidiacono2010beyond}. A department term of about 1\% of 4-year pay per
SD of department prestige therefore implies that department standing carries little information about
graduates' market productivity (small $\kappa_f\sigma_y$). The alternative, that employers cannot see
departments, fits only if learning about individuals is slow over the horizons observed ($b\approx0$ and
$k\gtrsim20$), against that evidence. The results do not separate information from content. Over careers,
learning moves the department loading away from zero if and only if $\kappa_f\neq0$, $\omega_f<1$ and $b<1$; it
raises the loading only if, in addition, $\kappa_f>0$.

\paragraph{Proposition 4 (experience).} (a) If employers' posterior mean for a unit equals the unit's mean,
learning about individuals cannot move program-mean pay: the program mean of each graduate's expected
productivity equals the program's mean productivity at every $t$. This is the result of \citet{farber1996learning}
that the return to a variable employers see at hiring does not change with experience, applied to program means.
The loading on a status measure net of selectivity then falls with experience, the signature of
\citet{altonji2001employer}, only if employers' beliefs about names are imprecise ($\omega_i<1$) and their entry
prior has a status component that the measure tracks beyond selectivity. (b) A rise in the institution-level
covariance of pay with status, in log points, requires at least one of four sources:
\begin{itemize}
\item growth correlated with status, such as access to firms with internal ladders and campus recruiting
\citep{gibbons1999theory,rivera2015pedigree,weinstein2018employer}; college reputation correlates with earnings
growth in \citet{macleod2017big};
\item learning about names that employers rarely see ($\omega_i<1$), which raises the covariance only if
employers' prior under-weights status (with a halo, learning lowers it);
\item migration of graduates of high-status institutions toward high-wage labour markets;
\item a change in who is observed, for example as graduate students enter the earnings data.
\end{itemize}
On standardised scales a loading also rises when the variance of pay unrelated to status shrinks. Fast learning
implies a concave rise, and a near-linear rise argues against fast learning. Experience profiles are concave,
however, so status-correlated growth can be concave too, and the shape of the rise does not separate learning
from growth. Where the rise occurs can (test T3).

\paragraph{Proposition 5 (pay schedules).} Suppose a share $\pi_f$ of a field's graduates is paid on a schedule,
and let $\bar\tau_{if}$ be the schedule level where program $(i,f)$'s graduates work. Treating the median as the
mean of log pay, to first order $\mathrm{Cov}(F,m)=(1-\bar\pi_f)\,\mathrm{Cov}_{\text{market}}(F,m)+\bar\pi_f\,
\mathrm{Cov}(F,\bar\tau)$, plus a sorting term when higher-status programs send fewer graduates into scheduled
jobs. Under a national schedule the second term is zero at entry, and coupling tends to the sorting term as
$\bar\pi_f\to1$. Over careers it can grow. If promotion across grades (NHS bands, teachers' leadership ranges)
or exit from scheduled jobs depends on status, $\bar\tau$ comes to vary with status. Under local schedules,
coupling runs through where graduates work. In such fields the wage does not measure the market's valuation of
status, so weak coupling there is not evidence that employers ignore status. The model's content there is a
prediction about quantities, namely who is hired where and who is promoted, which test T6 checks.

\paragraph{How the results line up.} No result follows from the structure for every admissible parameter value.
Every result is accommodated after the fact (\SItheory), at different costs.
\begin{itemize}
\item \emph{Accommodated by parameter values inferred from these same results.} Field prestige is no better
than $G$. The identity sets the threshold: at the across-field means, $r_{FG}=0.79$, $c_G=0.45$ and $c_F=0.40$
give $0.16$ against $\mathrm{sr}_F=0.08$, which requires a small $\kappa_f\Lambda_{ft}$. Agreement concentrates
in institution-wide status and selectivity and leaves a small departmental remainder. The status loading does not
decay net of selectivity, which requires $\omega_i\approx1$ or no status component in the prior that $G$ tracks.
Coupling is low in UK medicine and nursing and flat over their careers, given a schedule share coded from outside
that cannot be separated from field type. The field map is collinear with selectivity pricing, under an added
one-factor assumption.
\item \emph{Accommodated only by added elements.} The rise over careers needs growth, migration, composition or a
scale effect. Net of $G$, the selectivity loading also rises (UK $+0.024$ per year). That fits learning about
names, selectivity-correlated growth or a common rise of standardised loadings, and together with the flat status
loading it is not evidence that employers know names precisely. UK teaching's rise ($+0.050$ per year) needs the
grade-progression term, which was added after the rise was seen.
\item \emph{In tension.} At year 1, field prestige carries the association and $G$ does not ($b_F=+0.189$
\ci{+0.060}{+0.309}, $b_G=-0.013$), and $b_F$ stays flat to year 10. Learning alone produces this for no
parameter values: $\kappa_f\approx0$ contradicts $b_F>0$, $\kappa_f>0$ with $\omega_f<1$ makes $b_F$ rise, and
$\omega_f=1$ makes the name loading flat as well. Four readings remain: employers know departments from entry;
$\lambda_t$ barely changes over these horizons; $F$ tracks a trait other than the department; or year 1 is an
artefact, since graduate study and falling coverage depress year-1 medians at high-status institutions. Because of
the last, pay cannot be read as employer reputation at entry on these data. Coupling in US special education and
UK allied health despite pay schedules is also in tension.
\item \emph{Outside the model.} The Great Recession entry window, cross-window drift, the elite tail and
disciplinary kinship.
\end{itemize}

\paragraph{Registered tests.} Construct readings earn their names through predictions that could fail
\citep{cronbach1955construct}. We specified nine tests on 25 September 2026, before computing any statistic named
in them. Several reuse data analysed for other results, and T3 was chosen knowing that the $G$ loading rises.
The date is self-attested: the digest of the specification detects later edits, not when it was written
(\SItheory).
\begin{itemize}
\item \emph{Construct audits of $F$.} T0 asks whether coupling survives when $F$ is replaced by ranks corrected
for production and hiring volume. T1 asks whether one index orders departments both by where their PhDs are
hired and by whom they hire, benchmarked against what SpringRank's own model produces. T2 asks whether, among
the same bachelor's graduates, the department-specific part of $F$ predicts the rank of the PhD program they
enter (from ORCID records), against pay, for which the answer is already known; it is a consistency check with one
new prediction, not a test of the mechanism. A world of network position or placement power predicts T1 and T2 as
a world of valuation does; only T0 separates placement power.
\item \emph{Dynamics.} T3 asks whether the rise of the $G$ loading concentrates where employers rarely see an
institution's graduates (learning) or not (growth). T4 asks whether, among UK graduates with the same prior
attainment, the premium for a more selective institution falls with years since graduation, as employer
learning predicts if attainment is unseen.
\item \emph{Scope.} T5 asks whether the field map is the same with non-academic status indices (faculty
salaries, instructional spending, endowment). T6 asks whether, in teaching and nursing, status predicts
placement into regions with higher scheduled pay. T7 asks how much of coupling and of its rise survives the
wage level of where graduates work, and T8 how much of coupling runs between sectors.
\end{itemize}

\paragraph{What the model lets us claim.} \emph{Relational claims} need only classical measurement error. The
academy's hiring order and the market's pay order of the same institutions agree to degree $\rho_f$; the
agreement lies mostly at the institution's name, grows with years since graduation, and is weak where pay follows
a schedule. The Spearman-based statements survive any monotone rescaling of either indicator; the
within-institution estimates are in log pay and do not. \emph{Construct claims} add Proposition~1. Where the
academy's and the market's orderings agree, they agree at the institution's name, mostly in what both share with
the selectivity of its intake; below the name they agree little, clearly only in computer science. Within
institutions, one SD of department prestige goes with about 1\% of pay, and up to about 3\% under the lower-bound
reliability of $F$, so the department-specific part is diluted, not shown to be absent. These claims are not
netted of destination location, and they read $F$ as the academy's valuation only under vertical sorting net of
volume (T0, T1). \emph{Mechanism claims} are not made. Little agreement below the name may mean that department
standing carries little information about graduates' market productivity, or that employers cannot attribute
performance to departments and learn slowly about individuals. Evidence on the speed of employer learning
favours the first, and no registered test separates them. A failed T1 would withdraw only the reading of $F$ as
a belief about the quality of a department's PhDs, leaving $F$ as the ordering revealed by net placement.

\paragraph{What the model does not claim.}
\begin{itemize}
\item \emph{Causal value added.} The components $x$ and $y$ bundle selection, peers and value added, and every
proposition concerns covariances. Nothing here says what attending an institution or department does to pay.
\item \emph{That $F$ measures teaching quality, or research quality as such.} $F$ is the academy's revealed
ordering of departments as producers of researchers, it may be partly self-reinforcing status
\citep{burris2004academic}, and it may partly reflect production volume.
\item \emph{That wages equal productivity.} Competitive pricing is an assumption. Monopsony
\citep{staiger2010there}, rents, compensating differentials and pay schedules break it, and later-horizon pay
mixes reputation, learned productivity, growth and location.
\item \emph{That employers read academic rankings.} The name term can be a pure common cause ($\chi$), such as
resources that attract both researchers and able students.
\end{itemize}
